% ===================================================================
%  TAFEL Nr. 15 — Der Markt. --- Die Lebensmittel.
%  Traduction 2026 d'apres texto/io, controlee sur texto/fr et
%  texto/en. Consigne : voir texto/de/10-tafel-01.tex.
%
%  LE MARCHE EST L'ENDROIT OU LA LANGUE ALLEMANDE SE DIVISE, ET CETTE
%  PAGE EN FAIT LA DEMONSTRATION EN QUARANTE MOTS.
%
%  Partout ailleurs dans les seize tableaux, l'allemand du nord et
%  celui du sud disent la meme chose. Ici, non :
%     nord              sud et Autriche      la chose
%     Kartoffel         Erdapfel             la pomme de terre
%     Tomate            Paradeiser           la tomate
%     Aprikose          Marille              l'abricot
%     Pflaume           Zwetschke            la prune
%     Brötchen          Semmel               le petit pain
%     Sahne             Obers, Rahm          la creme
%     Blumenkohl        Karfiol              le chou-fleur
%     Grüne Bohne       Fisole               le haricot vert
%     Aubergine         Melanzani            l'aubergine
%     Hackfleisch       Faschiertes          la viande hachee
%     Johannisbeere     Ribisel              la groseille
%     Meerrettich       Kren                 le raifort
%  Douze couples sur une seule planche, et pas un ailleurs. LA
%  NOURRITURE EST CE QU'ON ACHETE TOUS LES JOURS ET CE QU'ON NE
%  DISCUTE JAMAIS AVEC UN ETRANGER : rien ne pousse le marche a
%  s'unifier, et rien ne l'a unifie.
%
%  CETTE COLONNE-CI ECRIT LA COLONNE DE GAUCHE, l'allemand standard
%  d'Allemagne, comme aux tableaux 5 et 12. C'est la troisieme fois
%  que le partage se presente et la premiere ou il porte sur douze
%  mots au lieu de deux. La colonne autrichienne, quand elle viendra,
%  aura ici sa page la plus differente des seize.
%
%  ET « ERDAPFEL » N'EST PAS UN PROVINCIALISME : c'est l'ancien nom
%  commun, la pomme-de-terre mot pour mot, celui que le francais a
%  garde. Le nord l'a remplace par « Kartoffel », qui est l'italien
%  « tartufolo », la TRUFFE — parce que le tubercule y ressemblait. Et
%  une troisieme forme, « Grundbirne », la poire-de-terre, est passee
%  chez les voisins slaves : krompir en serbe, krumpir en croate,
%  krumpli en hongrois. TROIS NOMS ALLEMANDS POUR UNE SEULE PLANTE,
%  ET C'EST LE TROISIEME, QUE PLUS PERSONNE NE DIT EN ALLEMAND, QUI A
%  VOYAGE LE PLUS LOIN.
%
%  LE TABLEAU 15 ESTONIEN AVAIT TROUVE LA MEME PLANTE ET LA MEME
%  HESITATION : « kartul » contre « maaõun », la pomme-de-terre. Il
%  avait garde l'emprunt. L'allemand du sud a garde le compose,
%  l'allemand du nord l'emprunt, et l'estonien l'emprunt aussi — mais
%  l'emprunt qu'il a pris, c'est celui du nord, comme tout le reste.
%
%  ENFIN, CHAQUE ALIMENT DE CETTE PAGE PORTE LE NOM DU BATEAU QUI L'A
%  APPORTE. Ceux du Nouveau Monde sont entres chacun par une porte
%  differente : Kartoffel par l'Italie, Tomate et Mais et Schokolade
%  par l'Espagne, Ananas par le Portugal. Ceux de l'Ancien sont venus
%  de plus loin encore : Pfeffer du latin, Zucker et Kaffee de
%  l'arabe, Reis du grec, Tee du chinois par le neerlandais, Safran de
%  l'arabe. Et le Truthahn de l'alinea 6 porte, comme
%  l'anglais « turkey », le nom d'un pays ou il n'a jamais vecu.
% ===================================================================

%%K t15-apar-1 apar
\VUsaut{4.0mm}
\VUtitre{15.0pt}{60}{TAFEL Nr. 15}
\VUsaut{3.0mm}

%%K t15-tit-1 sub
\VUpk{11.4pt}{40}{Der Markt. --- Die Lebensmittel.}
\VUsaut{3.0mm}

%%K t15-01-1 p
1. --- Die meisten Händler, die uns mit den nötigen Lebensmitteln versorgen, sind
unter der \VUgras{Halle} \textsuperscript{(1)} der \VUgras{Markthalle} oder in den
Nachbarstraßen versammelt.

%%K t15-02-1 p
2. --- Der \VUgras{Schlachter} \textsuperscript{(2)}, der \VUgras{Schinken}
\textsuperscript{(3)}, \VUgras{Blutwurst} \textsuperscript{(4)}, \VUgras{Wurst}
\textsuperscript{(5)}, \VUgras{Kaldaunenwurst} \textsuperscript{(6)} und
\VUgras{Speck} \textsuperscript{(7)} verkauft, steht neben der
\VUgras{Butter- und Käsehändlerin} \textsuperscript{(8)}, deren Stand mit
rotrindigen \VUgras{holländischen Käsen} \textsuperscript{(9)},
\VUgras{Camemberts} \textsuperscript{(10)}, großen \VUgras{Greyerzer Laiben}
\textsuperscript{(11)} und frischen \VUgras{Butterstücken} \textsuperscript{(12)}
ausgelegt ist.

%%K t15-03-1 p
3. --- Vor ihr bietet eine \VUgras{Gemüsehändlerin} \textsuperscript{(13)} ihren
Kundinnen schöne \VUgras{Tomaten} \textsuperscript{(14)}, \VUgras{Blumenkohl}
\textsuperscript{(15)}, \VUgras{Salat} \textsuperscript{(16)}, \VUgras{Kohl}
\textsuperscript{(17)}, \VUgras{Kürbisse} \textsuperscript{(19)}, \VUgras{Melonen}
\textsuperscript{(18)} und sogar \VUgras{Pilze} \textsuperscript{(20)} an. Aber ein
\VUgras{Brauereikutscher}, der seinen \VUgras{Lieferwagen} \textsuperscript{(21)}
voller \VUgras{Bierfässer} \textsuperscript{(22)} genau vor ihrem Stand angehalten
hat, versperrt ihren Kundinnen den Weg. Ein Gemüsegärtner bringt ihr einen Korb
\VUgras{Artischocken} \textsuperscript{(23)}.

%%K t15-tit-2 sub
\VUsaut{4.0mm}
\VUpk{10.2pt}{40}{Verkaufen und kaufen.}
\VUsaut{3.0mm}

%%K t15-04-1 p
4. --- Auf dem Markt herrscht großer Betrieb, weil es Zeit für die
\VUgras{Versteigerung} \textsuperscript{(24)} ist. Die \VUgras{Hausfrauen}
\textsuperscript{(26)}, die hier einkaufen, bleiben unterwegs am Stand der
\VUgras{Kaldaunenhändlerin} \textsuperscript{(25)} stehen, um \VUgras{Kaldaunen} oder
einen \VUgras{Kalbskopf} zu kaufen.

%%K t15-05-1 p
5. --- Eine stämmige \VUgras{Köchin} \textsuperscript{(27)} feilscht mit der
\VUgras{Fischhändlerin} \textsuperscript{(28)} um einen sehr schönen
\VUgras{Thunfisch}; die Händlerin setzt ihre Preise an, wie es ihr passt. Sie hat
eine große Lieferung \VUgras{Aale}, \VUgras{Seezungen}, \VUgras{Forellen} und
\VUgras{Karpfen} bekommen. Ihr \VUgras{Kabeljau} und ihre \VUgras{Muscheln} sind
sehr frisch.

%%K t15-tit-3 sub
\VUsaut{4.0mm}
\VUpk{10.2pt}{40}{Billige Ware und teure.}
\VUsaut{3.0mm}

%%K t15-06-1 p
6. --- Die \VUgras{Geflügelhändlerin} \textsuperscript{(29)}, die neben ihr steht, ist
sehr ehrlich. Wenn sie einen \VUgras{Truthahn}, eine \VUgras{Gans}, eine
\VUgras{Ente} oder ein schlichtes \VUgras{Huhn} verkauft, verlangt sie nicht mehr als
den angemessenen Preis.

%%K t15-07-1 p
7. --- An der Ecke des Platzes, im ersten Stock, hat sich vor kurzem eine
\VUgras{Weißwarenhändlerin} \textsuperscript{(30)} niedergelassen, wo die Käufer immer
eine sehr gute Auswahl an \VUgras{Tischtüchern}, \VUgras{Servietten},
\VUgras{Schürzen} und \VUgras{Geschirrtüchern} finden. Auf derselben Etage hat ein
\VUgras{Messerschmied} \textsuperscript{(31)} seine Werkstatt eingerichtet. Er
schleift die \VUgras{Messer} der Marktfrauen und macht den Gärtnern
\VUgras{Gartenmesser} aus dem härtesten \VUgras{Stahl}.

%%K t15-tit-4 sub
\VUsaut{4.0mm}
\VUpk{10.2pt}{40}{Das Kolonialwarengeschäft.}
\VUsaut{3.0mm}

%%K t15-08-1 p
8. --- Unter seiner Werkstatt hat vor kurzem ein großes Lebensmittelgeschäft
eröffnet. Es ist nicht mehr der schlichte, bescheidene \VUgras{Kolonialwarenladen}
\textsuperscript{(32)} von früher, der nur Waren des täglichen Bedarfs führte ---
\VUgras{Tee} \textsuperscript{(33)}, \VUgras{Schokolade} \textsuperscript{(34)},
\VUgras{Hutzucker} \textsuperscript{(37)} und \VUgras{Seife} \textsuperscript{(38)}.
Hier wird alles verkauft: \VUgras{Kekse}, \VUgras{Eingemachtes}
\textsuperscript{(35)}, \VUgras{Liköre} \textsuperscript{(36)}, \VUgras{Rum},
\VUgras{Branntwein}, \VUgras{Kirschwasser}, \VUgras{Anisett} und
\VUgras{Curaçao}. Wild ist hier ebenso leicht zu bekommen --- \VUgras{Hase}
\textsuperscript{(39)}, \VUgras{Drosseln} \textsuperscript{(40)},
\VUgras{Rebhühner} \textsuperscript{(41)}, \VUgras{Wachteln}
\textsuperscript{(42)}, \VUgras{Wildenten} \textsuperscript{(43)} oder
\VUgras{Fasan} \textsuperscript{(44)} --- wie tropisches Obst, etwa \VUgras{Bananen}
\textsuperscript{(46)} und \VUgras{Ananas}.

%%K t15-09-1 p
9. --- Ein sehr gefälliger \VUgras{Verkäufer} \textsuperscript{(45)} des Ladens ist
bereit, alles herauszugeben, was man verlangt: \VUgras{Marmelade}
\textsuperscript{(47)}, von Johannisbeeren oder Quitten, \VUgras{Schmalz}
\textsuperscript{(48)}, \VUgras{Essiggurken} \textsuperscript{(49)} oder gesalzene
\VUgras{Sardinen} \textsuperscript{(50)}.

%%K t15-09-2 p
Das ist für die \VUgras{Kunden} \textsuperscript{(51)} sehr praktisch, die neben der
gewöhnlichsten Ware auch die seltensten Frühgemüse und das beste Obst finden: saftige
\VUgras{Birnen} \textsuperscript{(52)}, \VUgras{Pflaumen} \textsuperscript{(53)},
\VUgras{Trauben} \textsuperscript{(54)}, \VUgras{Pfirsiche} \textsuperscript{(55)},
\VUgras{Mandeln} \textsuperscript{(56)} und \VUgras{Walnüsse}
\textsuperscript{(57)}.

%%K t15-tit-5 sub
\VUsaut{4.0mm}
\VUpk{10.2pt}{40}{Die Kunden.}
\VUsaut{3.0mm}

%%K t15-10-1 p
10. --- Der \VUgras{Reis} \textsuperscript{(58)} und die \VUgras{Linsen}
\textsuperscript{(59)} stehen neben der Kiste mit geräucherten \VUgras{Heringen}
\textsuperscript{(60)}; die \VUgras{Kartoffel} \textsuperscript{(61)} und die
\VUgras{Bohne} \textsuperscript{(63)} sitzen neben den \VUgras{Äpfeln}
\textsuperscript{(62)}, \VUgras{Feigen} \textsuperscript{(64)},
\VUgras{Backpflaumen} \textsuperscript{(65)} und \VUgras{Haselnüssen}
\textsuperscript{(66)}. Dieser Laden verkauft auch frische \VUgras{Eier}
\textsuperscript{(67)} und \VUgras{Oliven} \textsuperscript{(68)}, und so füllen sich
die \VUgras{Körbe} \textsuperscript{(70)} und die \VUgras{Netztaschen}
\textsuperscript{(73)} sehr schnell.

%%K t15-11-1 p
11. --- Der \VUgras{Kaffee} \textsuperscript{(72)} dieses ausgezeichneten Hauses hat
sich bei den \VUgras{Dienstmädchen} \textsuperscript{(69)} und den Köchinnen einen
wohlverdienten Namen gemacht, weil der alte \VUgras{Kolonialwarenhändler}
\textsuperscript{(71)} ihn mit größter Sorgfalt selbst röstet.

%%K t15-tit-6 sub
\VUsaut{4.0mm}
\VUpk{10.2pt}{40}{Was es zu sehen gibt.}
\VUsaut{3.0mm}

%%K t15-12-1 p
12. --- Nicht alle kommen zum Markt, um Lebensmittel zu kaufen. Im malerischen Hin
und Her der kleinen Gasse, die zur Halle führt, sieht man die verschiedensten Leute.
Neben einem gutmütigen \VUgras{Schlachthausarbeiter} \textsuperscript{(75)}, der eine
\VUgras{Schubkarre} \textsuperscript{(74)} vor sich herschiebt, stehen zwei Soldaten
auf Urlaub, sehr stolz auf ihre bunten Uniformen: der \VUgras{Husar}
\textsuperscript{(76)} in seinem blauen \VUgras{Dolman} und der \VUgras{Federtschako},
und der \VUgras{Zuavenkorporal} in weiten Pluderhosen mit einem \VUgras{Fes} auf dem
Kopf.

%%K t15-13-1 p
13. --- Der \VUgras{Bäcker} \textsuperscript{(80)}, der auf der Schwelle der
\VUgras{Bäckerei} \textsuperscript{(78)} steht, hat gerade den Teig für sein
\VUgras{Brot} \textsuperscript{(79)} geknetet und die \VUgras{Hörnchen}
\textsuperscript{(81)}, die \VUgras{Brötchen} \textsuperscript{(82)} und die
\VUgras{runden Laibe} \textsuperscript{(83)} aus dem Ofen geholt, die im Fenster
liegen.

%%K t15-14-1 p
14. --- Über der Bäckerei wohnt eine geschickte \VUgras{Schneiderin}
\textsuperscript{(84)}, die mehrere \VUgras{Stickerinnen} \textsuperscript{(85)}
beschäftigt. Neben ihr wohnt eine \VUgras{Wäscherin und Plätterin}
\textsuperscript{(86)}, die die \VUgras{Wäsche} \textsuperscript{(88)} stärkt, nachdem
sie sie gewaschen und getrocknet hat, und sie mit einem \VUgras{Bügeleisen}
\textsuperscript{(87)} plättet.

%%K t15-tit-7 sub
\VUsaut{4.0mm}
\VUpk{10.2pt}{40}{Fleisch, Fisch und Gemüse.}
\VUsaut{3.0mm}

%%K t15-15-1 p
15. --- In der \VUgras{Metzgerei} \textsuperscript{(89)} im Erdgeschoss wird Fleisch
jeder Art verkauft --- \VUgras{Hammelfleisch} \textsuperscript{(90)},
\VUgras{Rindfleisch} \textsuperscript{(94)} oder \VUgras{Kalbfleisch}
\textsuperscript{(92)} --- als \VUgras{Hammelkeulen}, \VUgras{Steaks},
\VUgras{Braten} und \VUgras{Koteletts}. Der \VUgras{Metzgergeselle}
\textsuperscript{(93)} zerlegt all dieses Fleisch und legt die
\VUgras{Markknochen} \textsuperscript{(96)} beiseite. An der Tür der Metzgerei steht
eine \VUgras{Austernhändlerin} \textsuperscript{(97)}, die \VUgras{Austern}
\textsuperscript{(98)} verkauft. Sie macht sie auf, wenn man will.

%%K t15-16-1 p
16. --- Alle diese Händler zahlen eine Gewerbeerlaubnis oder eine Standgebühr;
deshalb kennt der \VUgras{Polizist} \textsuperscript{(100)} kein Erbarmen mit den
armen \VUgras{fliegenden Gemüsehändlerinnen} \textsuperscript{(99)}, die ihnen
Konkurrenz machen, indem sie auf ihren Karren \VUgras{Karotten},
\VUgras{Radieschen}, \VUgras{Rüben}, \VUgras{Lauch} oder Bündel
\VUgras{Spargel}, \VUgras{frische Erbsen}, \VUgras{Spinat}, \VUgras{Sauerampfer},
\VUgras{Brunnenkresse} und \VUgras{Petersilie} herumfahren.

%%K t15-tit-8 sub
\VUsaut{4.0mm}
\VUpk{10.2pt}{40}{Achtung, der Polizist!}
\VUsaut{3.0mm}

%%K t15-17-1 p
17. --- Der Junge, der \VUgras{Knoblauch} \textsuperscript{(101)} und
\VUgras{Zwiebeln} verkauft, weiß sehr wohl, dass auch er seine Ware nicht in der
Nähe des Marktes verkaufen darf. Er läuft davon, auf die Gefahr hin, den Korb mit den
\VUgras{Krabben}, \VUgras{Flusskrebsen} und \VUgras{Garnelen} der
\VUgras{Händlerin} \textsuperscript{(102)} umzuwerfen, die \VUgras{Langusten} und
\VUgras{Hummer} verkauft.

%%K t15-18-1 p
18. --- Alle sind in Bewegung und machen großen Lärm; aber man hört trotzdem deutlich
den Ruf des wandernden \VUgras{Glasers} \textsuperscript{(103)} und den Schrei des
\VUgras{Kohlenhändlers} \textsuperscript{(104)}, der seinen Wagen eben einen
Augenblick stehen gelassen hat, um einen \VUgras{Sack Kohlen}
\textsuperscript{(105)} auszuliefern.
\VUsaut{6.0mm}
\VUfilet{22mm}
